\section{Method: Asynchronous Slide Evidence for Streaming Translation}
\label{sec:method}

\subsection{Design principle: off the critical path}

The central architectural commitment is that visual processing must never block
the audio decision loop. The streaming translator runs its READ/WRITE policy at
audio rate; stable slides usually change on a much slower timescale, whose
distribution we measure rather than assume. We therefore run a
\emph{slide worker} as an asynchronous process: it watches the video stream,
detects slide changes, reads the new slide once, structures the result into a
persistent evidence state, and atomically updates a shared \emph{evidence
cache}. The translator performs a small retrieval from this cache at each WRITE;
we measure this on-path cost rather than declaring it zero. This contrasts with
synchronous multimodal fusion
\citep{koneru2025omnifusion}, where image tokens ride the decoding path and
demonstrably add latency. In the worst case (the speaker names a term in the
seconds before the worker finishes a new slide), the system degrades exactly to
the audio-only baseline; visual evidence can arrive late, but can never slow
the audio path down.

\begin{figure}[t]
\centering
\fbox{\parbox{0.92\linewidth}{
\centering
\vspace{1.2em}
\textbf{audio path (ms-scale):} speech chunks $\rightarrow$ streaming
translator (READ/WRITE) $\rightarrow$ committed target words \\[0.6em]
\textbf{slide path (dwell-scale, async):} video frames $\rightarrow$
change detector $\rightarrow$ one-shot visual reader $\rightarrow$ semantic
evidence state $\rightarrow$ persistent cache $\uparrow$ retrieved at WRITE
\vspace{1.2em}
}}
\caption{Two-timescale architecture. The slide path updates a persistent
evidence cache once per change; the non-blocking audio path only retrieves from
it. Placeholder for final figure.}
\label{fig:overview}
\end{figure}

\subsection{Slide worker: what to extract}

\paragraph{Change detection.} Frame differencing on a low-rate sample
(1 frame / 2--5 s) with a stability window, so partial animations and camera
cuts to the speaker do not trigger re-reads; the current slide persists in the
store until replaced (lookback bounded at 90 s in our probe).

\paragraph{Structured visual reading with an OCR control.} The reader extracts
visible text together with formulae, layout, chart relations, emphasis, and
source regions. Strong OCR is a mandatory baseline, not a strawman. Earlier
text-prompt probes found that a VLM repaired fragmented 480p OCR and helped a
selected hard stratum, but they did not establish raw-pixel necessity. We claim
beyond-OCR value only when structured evidence beats strong OCR on preregistered
layout, formula, chart, or emphasis slices; direct image embeddings earn a
separate claim only if they beat both structured and oracle text-equivalent
evidence.

\paragraph{Evidence structuring.} We deduplicate content across consecutive
slides, language-tag terms, preserve region and relation provenance, and create
compact views for recognition support (M2) and target-form supply (M3).
Domain-specificity filtering against a general frequency list follows
\citet{sinhamahapatra2025slides}. The complete cached state remains available
for auditing even when the translator receives only a small packet.

\subsection{Need-gated injection: when to use it}
\label{sec:gating}

The probe's sharpest lesson is that \emph{what} is not enough: ungated VLM
terms are flat pooled ($-0.6$ chrF) because easy segments are slightly hurt by
irrelevant additions, cancelling the large hard-segment gains. Injecting only
where needed changes everything: gating on the (oracle) hard stratum yields
$+7.3$ chrF pooled ($p{<}0.001$), matching the oracle-terms upper bound. The
open problem is therefore \emph{need prediction}: deciding, online, that the
upcoming content is likely to defeat the audio-only translator.

We formalize the gate as a scoring rule over the current state,
\begin{equation}
  g(t) \;=\; \alpha\,u(t) \;+\; \beta\,\mathrm{match}(G_t, \hat{x}_{\le t})
  \;-\; \gamma\,\mathrm{staleness}(G_t),
  \label{eq:gate}
\end{equation}
where $u(t)$ is translator uncertainty (instability of consecutive hypotheses
--- LCP stall rate --- or token-level entropy when logits are available),
$\mathrm{match}$ measures phonetic/lexical affinity between stored slide terms
$G_t$ and the recent source hypothesis (pinyin match for zh, subword overlap
and loanword detection for European languages), and staleness discounts slides
that have been superseded. Evidence is injected at WRITE $t$ iff
$g(t) > \tau$, with terms rendered as a structured hint block (term, language,
source slide, confidence), never as raw slide text. A naive alternative ---
asking the translator LLM ``are these terms relevant?'' --- fails for lack of
selectivity (it fires on 88\% of segments, including 94\% of easy ones); the
gate must be grounded in \emph{failure prediction}, not topical relevance,
because slide terms are almost always topically relevant. Signal design also
matters: a retrospective matcher (fire when slide terms appear in the model's
own partial output) is provably too late --- by the time the term surfaces, no
help is needed --- and conservative instability thresholds under-fire on strong
translators (0/40 activations in our first implementation). Effective gates
must be \emph{forward-looking}: source-side lexical cues of incoming
terminology, phonetic affinity between the source stream and stored slide
terms, and sensitive (not consecutive) instability counters.

\subsection{Faithfulness constraints}

The injected hints carry an explicit instruction boundary: translate only what
has been spoken; hints may disambiguate or spell a spoken term but never add
visible-but-unspoken content. We evaluate this directly with wrong-slide and
wrong-time controls (adoption rate of planted bad evidence) and a
visible-but-unspoken hallucination check against the slide text. At probe
scale, prompt-level wrong terms are benign for a 32B translator (within
$\pm1$ chrF of baseline); we expect --- and test for --- larger risks under
image-input and fine-tuned injection modes. An injection-form study
(uniform and trie-constrained logit biasing vs.\ prompt injection) resolves
the apparent tension: logits-level biasing is benign but powerless (term
failures are phrase-planning failures, not single-token near-misses), while
the easy-segment harm once attributed to prompt injection traces to
\emph{serving nondeterminism} destabilizing Local Agreement commits --- under
deterministic decoding, prompt injection is simultaneously powerful ($+6.5$
chrF on hard segments, $p{<}0.001$) and benign (easy segments unharmed). The
full system therefore uses deterministic decoding (or a drift-tolerant commit
rule) with prompt-level hint injection; need-gating becomes a token-budget
optimization rather than a necessity.

\subsection{Training-free core, learned extensions}

The core pipeline is training-free (VLM reader + rule-based gate thresholds
tuned on development talks), which keeps it portable across base translators.
Learned extensions fit two places without touching the architecture: a small
need-prediction classifier trained on (state, baseline-failure) pairs mined
automatically from development runs (our pilot: 1{,}540 auto-labeled segments;
draft-translation logprob features reach talk-held-out AUC 0.62 where source
surface features are at chance, preserving the full hard-stratum gain
end-to-end while easy-segment false fires remain the cost), and preference-tuning the translator to follow hint
blocks conservatively (use terms when spoken, ignore otherwise) using
synthetic slide augmentation \citep{sinhamahapatra2025slides} --- synthetic
data is admissible in training precisely because our test sets forbid it.
